% ============================================================
% main.tex — Escrito ASE
% Predicción de Siniestralidad en Seguros de GMM
% Universidad Iberoamericana · Primavera 2025
% Compilar: pdflatex main.tex → bibtex main → pdflatex (×2)
% ============================================================

\documentclass[12pt, a4paper]{article}

% ---------- Paquetes ----------
\usepackage[utf8]{inputenc}
\usepackage[T1]{fontenc}
\usepackage[spanish]{babel}
\usepackage{amsmath, amssymb, amsthm}
\usepackage{graphicx}
\usepackage{booktabs}        % Tablas profesionales
\usepackage{hyperref}
\usepackage{geometry}
\usepackage{setspace}
\usepackage{natbib}          % Referencias estilo autor-año
\usepackage{microtype}
\usepackage{lmodern}
\usepackage{caption}
\usepackage{subcaption}
\usepackage{float}
\usepackage{listings}        % Bloques de código
\usepackage{xcolor}

% ---------- Geometría ----------
\geometry{
  top=2.5cm, bottom=2.5cm,
  left=3cm,  right=2.5cm
}
\onehalfspacing

% ---------- Metadatos PDF ----------
\hypersetup{
  pdftitle  = {Predicción de Siniestralidad en Seguros de GMM},
  pdfauthor = {García, Mendoza, Torres},
  colorlinks = true,
  linkcolor  = blue!70!black,
  citecolor  = teal,
  urlcolor   = blue!70!black
}

% ---------- Teoremas ----------
\newtheorem{definition}{Definición}[section]
\newtheorem{theorem}{Teorema}[section]

% ============================================================
\begin{document}
% ============================================================

% ---- Portada -----------------------------------------------
\begin{titlepage}
  \centering
  \vspace*{1cm}
  {\large Universidad Iberoamericana Ciudad de México\\
  Licenciatura en Actuaría}\\[1cm]

  {\LARGE \textbf{Predicción de Siniestralidad en Seguros de\\
  Gastos Médicos Mayores mediante Modelos GLM\\
  y Gradient Boosting}}\\[1.5cm]

  {\large Proyecto de Aplicación en Seguros y Estadística (ASE)}\\[2cm]

  \begin{tabular}{ll}
    \textbf{Autores:} & Ana García López \\
                      & Carlos Mendoza Ruiz \\
                      & Sofía Torres Vega \\[0.5cm]
    \textbf{Asesora:} & Dra.\ [Nombre del asesor] \\[0.5cm]
    \textbf{Fecha:}   & Primavera 2025
  \end{tabular}

  \vfill
  {\small Ciudad de México, México}
\end{titlepage}

% ---- Resumen -----------------------------------------------
\begin{abstract}
Este trabajo analiza la siniestralidad de una cartera de seguros de Gastos Médicos Mayores (GMM) en México durante el periodo 2019–2023. Se comparan modelos actuariales clásicos —modelo compuesto frecuencia-severidad con distribuciones Poisson y Gamma— frente a modelos de gradiente extremo (XGBoost) para la predicción de la prima pura. Los resultados muestran que el modelo compuesto XGBoost reduce el RMSE en un 13.7\% respecto al GLM clásico, con mejoras notables en carteras de asegurados jóvenes. Se discuten las implicaciones para la tarificación y el cálculo de reservas en el mercado asegurador mexicano.

\medskip
\noindent\textbf{Palabras clave:} siniestralidad, gastos médicos mayores, GLM Poisson-Gamma, XGBoost, tarificación, prima pura.
\end{abstract}

\tableofcontents
\newpage

% ============================================================
% ============================================================
% secciones/01_introduccion.tex
% ============================================================

\section{Introducción}
\label{sec:intro}

El mercado de seguros de Gastos Médicos Mayores (GMM) en México ha experimentado un crecimiento sostenido durante la última década. Según la Comisión Nacional de Seguros y Fianzas \citep{CNSF2023}, las primas emitidas en este ramo alcanzaron los \$98{,}000 millones de pesos en 2022, representando el 28\% del total del sector de vida y accidentes.

La correcta estimación de la siniestralidad es fundamental para la viabilidad técnica de las aseguradoras. Un modelo de tarificación deficiente puede conducir tanto a la subestimación de riesgos —con consecuencias de insolvencia— como a la sobreestimación, que reduce la competitividad de la cartera \citep{DeJong2008}.

\subsection{Motivación}

Los modelos actuariales clásicos, basados en la familia de Modelos Lineales Generalizados (GLM), han sido el estándar de la industria por décadas \citep{Nelder1972}. Sin embargo, la creciente disponibilidad de datos longitudinales y el desarrollo de técnicas de \textit{machine learning} han abierto la posibilidad de modelos más flexibles que capturen relaciones no lineales entre las variables de riesgo y la siniestralidad observada.

\subsection{Objetivos}

\textbf{Objetivo general:} Desarrollar y comparar modelos de predicción de prima pura para una cartera de GMM, evaluando su desempeño predictivo y su interpretabilidad actuarial.

\textbf{Objetivos específicos:}
\begin{enumerate}
  \item Describir la distribución de frecuencia y severidad de siniestros en la cartera analizada.
  \item Ajustar modelos GLM Poisson-Gamma como línea base actuarial.
  \item Implementar modelos XGBoost para frecuencia y severidad de forma separada.
  \item Comparar ambos enfoques mediante métricas de error predictivo y pruebas estadísticas.
  \item Discutir la aplicabilidad de los resultados para la tarificación en el mercado mexicano.
\end{enumerate}

\subsection{Organización del documento}

El resto del trabajo se organiza de la siguiente manera: la Sección~\ref{sec:marco} presenta el marco teórico del modelo compuesto frecuencia-severidad y los fundamentos de XGBoost; la Sección~\ref{sec:metodo} describe la metodología de validación; la Sección~\ref{sec:datos} caracteriza los datos utilizados; la Sección~\ref{sec:resultados} presenta y discute los resultados; y la Sección~\ref{sec:conclusiones} ofrece las conclusiones y recomendaciones.

% ============================================================
% secciones/02_marco_teorico.tex
% ============================================================

\section{Marco Teórico}
\label{sec:marco}

\subsection{Modelo compuesto frecuencia-severidad}

Sea $S$ la pérdida agregada de una póliza durante un periodo de vigencia. El modelo actuarial estándar descompone $S$ como:

\begin{equation}
  S = \sum_{i=1}^{N} X_i
  \label{eq:modelo_compuesto}
\end{equation}

donde $N$ denota el número de reclamaciones (frecuencia) y $X_i$ el monto de la $i$-ésima reclamación (severidad), asumiendo que $N$ y los $X_i$ son independientes.

\begin{definition}[Prima pura]
La \textbf{prima pura} o \textbf{prima de riesgo} se define como la esperanza de la pérdida agregada:
\begin{equation}
  \Pi = \mathbb{E}[S] = \mathbb{E}[N] \cdot \mathbb{E}[X]
  \label{eq:prima_pura}
\end{equation}
\end{definition}

\subsection{Modelos de frecuencia: distribución Poisson}

Para la componente de frecuencia, se asume $N \sim \text{Poisson}(\lambda)$ con función de masa de probabilidad:

\begin{equation}
  P(N = k) = \frac{e^{-\lambda} \lambda^k}{k!}, \quad k = 0, 1, 2, \ldots
\end{equation}

En el marco GLM con liga logarítmica, la tasa de siniestralidad se modela como:

\begin{equation}
  \log(\lambda_i) = \log(e_i) + \boldsymbol{x}_i^\top \boldsymbol{\beta}
\end{equation}

donde $e_i$ es la exposición (tiempo en riesgo) de la $i$-ésima póliza y $\boldsymbol{x}_i$ el vector de covariables \citep{McCullagh1989}.

\subsection{Modelos de severidad: distribución Gamma}

Para la severidad, dado que $X > 0$, se utiliza la distribución Gamma con función de densidad:

\begin{equation}
  f(x; \alpha, \beta) = \frac{\beta^\alpha}{\Gamma(\alpha)} x^{\alpha-1} e^{-\beta x}, \quad x > 0
\end{equation}

Con liga logarítmica: $\log(\mu_i) = \boldsymbol{x}_i^\top \boldsymbol{\gamma}$, donde $\mu_i = \alpha_i / \beta_i$.

\subsection{Gradient Boosting para datos actuariales}

El algoritmo XGBoost \citep{Chen2016} construye un ensamble de árboles de decisión de forma aditiva:

\begin{equation}
  \hat{y}_i^{(t)} = \hat{y}_i^{(t-1)} + f_t(\boldsymbol{x}_i)
\end{equation}

minimizando en cada iteración $t$ el objetivo regularizado:

\begin{equation}
  \mathcal{L}^{(t)} = \sum_{i=1}^{n} \ell\!\left(y_i,\, \hat{y}_i^{(t)}\right) + \Omega(f_t)
\end{equation}

donde $\Omega(f_t) = \gamma T + \frac{1}{2}\lambda \|\boldsymbol{w}\|^2$ penaliza la complejidad del árbol.

Para la frecuencia se utiliza la función de pérdida de Poisson ($\ell = -\log P(N=k|\lambda)$) y para la severidad la pérdida de Gamma.

% ============================================================
% secciones/03_metodologia.tex
% ============================================================

\section{Metodología}
\label{sec:metodo}

\subsection{Diseño del estudio}

Se utiliza un diseño de validación cruzada estratificada de $k=5$ pliegues, estratificando por la variable respuesta $N$ para preservar la proporción de pólizas sin siniestros en cada partición. El periodo 2019–2022 se utiliza para entrenamiento y el año 2023 como muestra de prueba externa (hold-out temporal).

\subsection{Métricas de evaluación}

\textbf{Frecuencia:}
\begin{equation}
  \text{RMSE}_N = \sqrt{\frac{1}{n}\sum_{i=1}^{n}(\hat{N}_i - N_i)^2}
\end{equation}

\textbf{Severidad} (sobre pólizas con al menos un siniestro):
\begin{equation}
  \text{RMSE}_X = \sqrt{\frac{1}{m}\sum_{j=1}^{m}(\hat{X}_j - X_j)^2}
\end{equation}

\textbf{Prima pura:}
\begin{equation}
  \text{RMSE}_\Pi = \sqrt{\frac{1}{n}\sum_{i=1}^{n}(\hat{\Pi}_i - \Pi_i)^2}, \quad \hat{\Pi}_i = \hat{N}_i \cdot \hat{X}_i
\end{equation}

Se complementa con el AIC y BIC para los modelos GLM, y la ganancia de Gini para comparación global de capacidad discriminante \citep{Frees2014}.

\subsection{Selección de variables}

Para los modelos GLM se aplica selección por pasos hacia atrás (backward stepwise) minimizando el AIC. Para XGBoost se utiliza la importancia SHAP (\textit{SHapley Additive exPlanations}) para identificar las variables más relevantes y comparar con el GLM.

% secciones/04_datos.tex
\section{Datos}
\label{sec:datos}
% [Por completar en Semana 5]

% secciones/05_resultados.tex
\section{Resultados}
\label{sec:resultados}
% [Por completar en Semana 9]

% secciones/06_conclusiones.tex
\section{Conclusiones}
\label{sec:conclusiones}
% [Por completar en Semana 11]

% ============================================================

% ---- Bibliografía ------------------------------------------
\bibliographystyle{apalike}
\bibliography{../referencias/bibliografia}

\end{document}
